\documentclass[11pt]{article}
\usepackage[margin=1in]{geometry}
\usepackage{amsmath,amssymb,amsthm}
\usepackage{graphicx}
\usepackage{microtype}
\usepackage[hidelinks]{hyperref}

\newtheorem{theorem}{Theorem}
\newtheorem{lemma}[theorem]{Lemma}
\newtheorem{corollary}[theorem]{Corollary}
\newcommand{\MIN}{\mathrm{MIN}}
\newcommand{\MAX}{\mathrm{MAX}}
\newcommand{\cZ}{\mathcal Z}
\newcommand{\M}{\mathrm M}
\newcommand{\oplusone}{\mathop{\oplus}\limits_{1}}

\title{No almost complete embedding for the original\\
Kalton operator-space pair of arXiv:2106.15022}
\author{Automated mathematics run \texttt{fa\_banach\_001}}
\date{19 July 2026}

\begin{document}
\maketitle

\noindent\textbf{Status.} Claimed full solution, likely valid, pending human
review.  The proof also gives a stronger obstruction for arbitrary
$\ell_1$-sums of minimal operator spaces.

\section{Source question}

Braga's Theorem 1.8 constructs operator spaces $X$ and $Y$ such that $X$
embeds linearly isomorphically, but not completely isomorphically, into $Y$,
while the bounded subsets of $X$ and $Y$ are almost completely coarsely
equivalent \cite{Braga2021}.  The paragraph immediately following the theorem
asks whether the \emph{specific} spaces constructed there also satisfy that
$X$ does not almost completely isomorphically embed into $Y$.

\begin{figure}[ht]
  \centering
  \includegraphics[width=\textwidth]{figures/open_problem_crop.png}
  \caption{Theorem 1.8 and the source question on PDF page 4 of
  \cite{Braga2021}.}
\end{figure}

For the explicit construction, put $H=R\cap C$.  Let
$Q:L_\infty[0,1]\to H$ be the complete quotient adjoint to the standard
complete embedding $R+C\to L_1[0,1]$.  On the vector space
$L_\infty[0,1]$, let
\[
 \|y\|_{Y_j}=\max\{2^{-j}\|y\|_\infty,\|Qy\|_H\},
 \qquad
 \cZ(Q)=\Big(\bigoplus_{j\geq1}Y_j\Big)_{\ell_1}.
\]
The source pair is
\[
 Y=\cZ(Q),\qquad X=H\oplus\ker\widetilde Q.
\]
An almost complete isomorphic embedding of an operator space $E$ into $F$
means a sequence $T_n:E\to F$ and constants $a,b>0$ such that
\[
 a\|z\|_{\M_n(E)}\leq\|(T_n)_n z\|_{\M_n(F)}
 \leq b\|z\|_{\M_n(E)}
 \tag{1}\label{eq:almost}
\]
for every $n$ and $z\in\M_n(E)$.

\section{Proof intuition}

There are two steps.  First, the quotient coordinates of every bounded map
$H\to\cZ(Q)$ form a compact operator into $\ell_1(H)$.  Therefore, on a
suitable infinite-dimensional Hilbert subspace, those coordinates are as
small as desired.  This reduces a hypothetical almost complete embedding into
$\cZ(Q)$ to one into
\[
 W=\Big(\bigoplus_{j\geq1}\MIN(L_\infty[0,1])\Big)_{\ell_1}.
\]

Second, fix one map $T_n:H\to W$.  Weak compactness gives a finite coordinate
head whose scalar-norm tail is tiny.  On that finite head, concentration on
high-dimensional Euclidean spheres supplies an $n$-dimensional Hilbert
subspace $F$ on which all coordinate norms are nearly constant
simultaneously.  Their operator norms on $F$ then have bounded sum.  For an
orthonormal basis $x_1,\ldots,x_n$ of $F$, the row matrix
$z=\sum_k e_{1k}\otimes x_k$ has norm $\sqrt n$ in $R\cap C$, whereas its
image in the finite minimal head has uniformly bounded norm.  The tiny tail
stays bounded after amplification.  This contradicts the lower estimate in
\eqref{eq:almost}.

\section{Three auxiliary lemmas}

\begin{lemma}[Uniform tails]\label{lem:tails}
Let $(E_j)$ be Banach spaces and
$Z=(\bigoplus_jE_j)_{\ell_1}$.  Every relatively weakly compact subset
$K\subset Z$ has uniformly small coordinate tails.  Consequently, if $G$ is
reflexive and $T:G\to Z$ is bounded, then
\[
 \|(I-P_N)T\|\longrightarrow0,
\]
where $P_N$ is the projection onto the first $N$ coordinates.
\end{lemma}

\begin{proof}
Suppose uniform tails fail.  A gliding-hump selection gives $\varepsilon>0$,
pairwise successive finite intervals $I_m\subset\mathbb N$, and $z^{(m)}\in K$
such that
\[
 \sum_{j\in I_m}\|z^{(m)}_j\|>\varepsilon.
\]
Choose $f_{m,j}\in B_{E_j^*}$, adjusting complex phases if necessary, so that
\[
 \left|\sum_{j\in I_m}f_{m,j}(z^{(m)}_j)\right|>\varepsilon.
\]
The map
\[
 A:Z\to\ell_1,
 \qquad
 (Az)_m=\sum_{j\in I_m}f_{m,j}(z_j),
\]
is contractive because the intervals are disjoint.  Thus $A(K)$ is relatively
weakly compact.  By the Schur property of $\ell_1$, it is relatively norm
compact and hence has uniformly small scalar coordinate tails.  This
contradicts $|(Az^{(m)})_m|>\varepsilon$.  The last assertion follows because
$T(B_G)$ is relatively weakly compact.
\end{proof}

\begin{lemma}[Simultaneous Hilbertian stabilization]\label{lem:stabilize}
Let $G$ be an infinite-dimensional Hilbert space and let
$p_1,\ldots,p_N$ be continuous seminorms on $G$ satisfying
\[
 \sum_{j=1}^Np_j(x)\leq C\|x\|_2\qquad(x\in G).
 \tag{2}\label{eq:seminormbound}
\]
For every $d\in\mathbb N$ and $\varepsilon>0$, there is a $d$-dimensional
subspace $F\subset G$ such that
\[
 \sum_{j=1}^N\sup_{x\in S_F}p_j(x)\leq C+\varepsilon.
 \tag{3}\label{eq:stabilized}
\]
\end{lemma}

\begin{proof}
We first recall the following additive form of Euclidean-sphere
stabilization.  Given $d$, $L$ and $\delta>0$, there is $D$ such that, whenever
$E$ is a Euclidean space of dimension at least $D$ and $p$ is a seminorm with
$p(x)\leq L\|x\|_2$, some $d$-dimensional $F\subset E$ satisfies
\[
 \sup_{S_F}p-\inf_{S_F}p<\delta. \tag{4}\label{eq:onep}
\]
Indeed, $p$ is $L$-Lipschitz on $S_E$.  L\'evy's concentration lemma says
that it is exponentially concentrated about a median.  For a random
isometric copy of $\ell_2^d$ in $E$, take a $\rho$-net of $S^{d-1}$ of
cardinality at most $(1+2/\rho)^d$.  If $L>0$, set
$t=\delta/4$ and $\rho=\delta/(4L)$ (decreasing $\rho$ to at most $1$ causes
no problem).  L\'evy's estimate and a union bound show, once $\dim E$ is large
enough in terms of $d,L,\delta$, that every net value differs from the median
by at most $t$.  Every sphere point then differs from the median by at most
$t+L\rho\leq\delta/2$, so the oscillation is at most $\delta$.  The case
$L=0$ is immediate.  This is the standard concentration proof of the additive
Dvoretzky lemma \cite[Chapter 2]{Ledoux}.

Choose positive $\delta_j$ with $\sum_j\delta_j<\varepsilon$.  Starting with
a sufficiently large finite-dimensional subspace of $G$, apply
\eqref{eq:onep} successively to $p_1,\ldots,p_N$, choosing the intermediate
dimensions backwards so that the final dimension is $d$.  Previously obtained
oscillation estimates persist on further subspaces.  On the resulting $F$,
put $a_j=\inf_{x\in S_F}p_j(x)$.  Then
\[
 \sup_{S_F}p_j\leq a_j+\delta_j.
\]
For any fixed $x\in S_F$, \eqref{eq:seminormbound} gives
$\sum_ja_j\leq\sum_jp_j(x)\leq C$.  Summing the preceding inequalities proves
\eqref{eq:stabilized}.
\end{proof}

\begin{lemma}[The minimal row identity]\label{lem:minrow}
Let $E$ be a Banach space, $S:F\to E$ a bounded map on an
$n$-dimensional Hilbert space, and $x_1,\ldots,x_n$ an orthonormal basis of
$F$.  Then
\[
 \left\|\sum_{k=1}^ne_{1k}\otimes Sx_k\right\|_{
 \M_n(\MIN(E))}=\|S\|.
 \tag{5}\label{eq:minrow}
\]
For the same orthonormal family regarded in $H=R\cap C$,
\[
 \left\|\sum_{k=1}^ne_{1k}\otimes x_k\right\|_{M_n(H)}=\sqrt n.
 \tag{6}\label{eq:rcrow}
\]
\end{lemma}

\begin{proof}
By the definition of the minimal operator-space structure, the left side of
\eqref{eq:minrow} is
\[
 \sup_{\phi\in B_{E^*}}
 \left(\sum_{k=1}^n|\phi(Sx_k)|^2\right)^{1/2}
 =\|(S|_F)^*\|=\|S\|.
\]
For \eqref{eq:rcrow}, unitary homogeneity lets us take $x_k$ to be the first
$n$ standard basis vectors of $\ell_2$.  The row component of the displayed
matrix has norm $\sqrt n$, while its column component has norm $1$.  The
$R\cap C$ norm is their maximum.
\end{proof}

\section{The minimal-sum obstruction}

\begin{theorem}\label{thm:minsum}
For every sequence of Banach spaces $(E_j)$, the operator space
\[
 W_E=\Big(\bigoplus_{j\geq1}\MIN(E_j)\Big)_{\ell_1}
\]
does not admit an almost complete isomorphic embedding of $R\cap C$.
\end{theorem}

\begin{proof}
Suppose that maps $T_n:H=R\cap C\to W_E$ satisfy \eqref{eq:almost} for fixed
$a,b>0$.  Fix $\eta>0$ and then a large $n$.  At scalar level,
$\|T_n\|\leq b$.  Since $H$ is reflexive, Lemma \ref{lem:tails} gives $N$ such
that
\[
 \|(I-P_N)T_n\|\leq\eta n^{-3/2}. \tag{7}\label{eq:tailchoice}
\]
Write $S_j:H\to E_j$ for the $j$th coordinate of $T_n$.  The underlying
Banach norm of $W_E$ is the $\ell_1$ norm, so
\[
 \sum_{j=1}^N\|S_jx\|\leq b\|x\|_2.
\]
Apply Lemma \ref{lem:stabilize} to $p_j(x)=\|S_jx\|$.  We obtain an
$n$-dimensional $F\subset H$ such that
\[
 \sum_{j=1}^N\|S_j|_F\|\leq b+\eta. \tag{8}\label{eq:headbound}
\]

Let $x_1,\ldots,x_n$ be an orthonormal basis of $F$ and set
\[
 z_F=\sum_{k=1}^ne_{1k}\otimes x_k\in\M_n(H).
\]
By Lemma \ref{lem:minrow}, $\|z_F\|=\sqrt n$.  The complete isometry of each
coordinate inclusion into the operator-space $\ell_1$ sum, the triangle
inequality, \eqref{eq:minrow}, and \eqref{eq:headbound} give
\[
 \|(P_NT_n)_n z_F\|
 \leq\sum_{j=1}^N\|(S_j)_n z_F\|
 =\sum_{j=1}^N\|S_j|_F\|
 \leq b+\eta. \tag{9}\label{eq:headmatrix}
\]

For every bounded map $A$ between operator spaces,
$\|A_n\|\leq n\|A\|$.  (Represent the target concretely and use that the norm
of an $n$ by $n$ operator matrix is at most $n$ times its largest entry
norm.)  Hence \eqref{eq:tailchoice} and \eqref{eq:rcrow} imply
\[
 \|((I-P_N)T_n)_n z_F\|
 \leq n\,\eta n^{-3/2}\sqrt n=\eta. \tag{10}\label{eq:tailmatrix}
\]
Combining \eqref{eq:headmatrix} and \eqref{eq:tailmatrix},
\[
 \|(T_n)_n z_F\|\leq b+2\eta.
\]
But the lower estimate in \eqref{eq:almost} gives
$\|(T_n)_n z_F\|\geq a\sqrt n$, a contradiction for large $n$.
\end{proof}

\section{Compact quotient coordinates and the original pair}

Set
\[
 W=\Big(\bigoplus_{j\geq1}\MIN(L_\infty[0,1])\Big)_{\ell_1},
 \qquad
 U=\Big(\bigoplus_{j\geq1}H\Big)_{\ell_1}.
\]
Here the canonical operator-space structure of the commutative
$C^*$-algebra $L_\infty[0,1]$ is its minimal operator-space structure.
For $y=(y_j)\in\cZ(Q)$, define
\[
 P_1y=(2^{-j}y_j)_j\in W,
 \qquad
 P_2y=(Qy_j)_j\in U.
\]
The map $\Gamma=(P_1,P_2):\cZ(Q)\to W\oplus_\infty U$ is a complete
isomorphic embedding with
\[
 \|\Gamma\|_{\rm cb}\leq1,
 \qquad
 \|\Gamma^{-1}:\Gamma(\cZ(Q))\to\cZ(Q)\|_{\rm cb}\leq2.
 \tag{11}\label{eq:gamma}
\]
Indeed, coordinatewise this is the defining maximum norm of $Y_j$.
Regrouping the two $\ell_1$ sums is completely contractive; the reverse map
is the sum of its two coordinate maps, each completely contractive, and so
has cb-norm at most $2$.

\begin{lemma}\label{lem:quotcompact}
For every bounded $u:H\to\cZ(Q)$, the quotient-coordinate operator
\[
 V_u:H\to U,\qquad V_ux=(Qu_jx)_j=P_2u(x),
\]
is compact.
\end{lemma}

\begin{proof}
First consider a bounded $T:H\to L_\infty[0,1]$.  We claim that $QT:H\to H$
is compact.  Otherwise, since its domain and range are Hilbert spaces, there
is an infinite-dimensional closed $E\subset H$ on which $A=QT$ is bounded
below.  The range $A(E)$ is closed and
\[
 \Pi=T|_E\,(A|_E)^{-1}P_{A(E)}Q:L_\infty[0,1]\to T(E)
\]
is a bounded projection onto a copy of $\ell_2$.  This contradicts the
classical complemented-subspace theorem for $L_\infty$ used in Braga's proof
of Proposition 4.6 \cite{Braga2021,AlbiacKalton}.  Thus each coordinate
$Qu_j$ is compact.

The map $V_u:H\to\ell_1(H)$ is bounded by the definition of $\cZ(Q)$.  Since
$H$ is reflexive, it is weakly compact.  Lemma \ref{lem:tails} shows that its
coordinate tails converge to zero in operator norm.  Its finite coordinate
heads are compact, so $V_u$ is a norm limit of compact operators.
\end{proof}

\begin{theorem}\label{thm:originalH}
The operator space $H=R\cap C$ does not almost completely isomorphically
embed into the original $\cZ(Q)$ of Theorem 1.8 in \cite{Braga2021}.
\end{theorem}

\begin{proof}
Suppose maps $u^{(n)}:H\to\cZ(Q)$ satisfy \eqref{eq:almost}.  Fix $n$.
By Lemma \ref{lem:quotcompact}, $P_2u^{(n)}$ is compact.  Approximate it in
operator norm by a finite-rank map and pass to the infinite-dimensional
kernel of that approximant.  We obtain an infinite-dimensional $F_n\subset H$
on which
\[
 \|P_2u^{(n)}|_{F_n}\|<\frac{a}{4n}.
\]
Every infinite-dimensional Hilbert subspace of $R\cap C$ is the range of a
complete isometry $J_n:H\to F_n$, because row and column norms are invariant
under Hilbert-space unitaries.

Put $T_n=P_1u^{(n)}J_n:H\to W$.  From \eqref{eq:almost} and
\eqref{eq:gamma},
\[
 \|\Gamma_nu^{(n)}_nJ_{n,n}z\|\geq\frac a2\|z\|.
\]
The $P_2$ component has $n$th amplification norm less than $a/4$.  Therefore
\[
 \|(T_n)_nz\|\geq\frac a2\|z\|,
 \qquad
 \|(T_n)_nz\|\leq b\|z\|
\]
for all $z\in\M_n(H)$.  Thus $(T_n)$ would almost completely embed $H$ into
$W$, contradicting Theorem \ref{thm:minsum}.
\end{proof}

\begin{corollary}[Answer to the source question]\label{cor:answer}
For the exact operator spaces
\[
 X=(R\cap C)\oplus\ker\widetilde Q,
 \qquad Y=\cZ(Q)
\]
constructed in Theorem 1.8 of \cite{Braga2021}, $X$ does not almost
completely isomorphically embed into $Y$.
\end{corollary}

\begin{proof}
The canonical $R\cap C$ summand of $X$ is completely isometric.  Restricting
any hypothetical almost complete embedding $X\to Y$ to this summand would
contradict Theorem \ref{thm:originalH}.
\end{proof}

\section{Verification notes and novelty scope}

The proof is qualitative; no computation is used.  The main points checked
independently in the proof are:
\begin{itemize}
\item Lemma \ref{lem:tails} is proved for arbitrary Banach-valued $\ell_1$
sums, so no hidden reflexivity of the coordinate spaces is assumed.
\item Lemma \ref{lem:stabilize} uses only finite families.  Intermediate
dimensions may be enormous, but the domain Hilbert space is infinite
dimensional and the final required dimension is exactly $n$.
\item The equality in Lemma \ref{lem:minrow} follows directly from the
definition of $\MIN(E)$; it does not rely on a special representation of
$L_\infty$.
\item The tail exponent $n^{-3/2}$ exactly compensates the general loss
$\|A_n\|\leq n\|A\|$ and the input row norm $\sqrt n$.
\item The quotient-coordinate reduction uses only compactness and the
complete distortion bound $2$ in (11); all constants have strict slack.
\end{itemize}

The required bounded novelty search was performed on 19 July 2026.  The four
local run indexes and the locally stored arXiv sources 2106.15022 and
2211.11854 were searched for the arXiv id, the exact question phrase,
``$R\cap C$'', ``almost complete isomorphic embedding'', ``minimal'', and
``$\ell_1$-sum''.  Bounded external searches used the same phrases, the title,
and the authors Braga and Oikhberg, restricted in the technical search to
primary arXiv records.  The search found the source paper, Braga--Oikhberg's
different-space existential strengthening \cite{BragaOikhberg2022}, and the
later small-scale nonlinear paper, but no proof for the exact original
$L_\infty\to R\cap C$ pair and no statement of Theorem \ref{thm:minsum}.
This is evidence only, not a guarantee of novelty.

Braga--Oikhberg say that their existential theorem addresses the motivation
of the question after Theorem 1.8, but their construction uses the different
quotient $\MAX(L_1)\to\MIN(\ell_2)$ and different spaces.  It does not settle
the exact-space assertion in Corollary \ref{cor:answer}.

\medskip
\noindent\textbf{Human-review recommendation.}  Review first the simultaneous
stabilization lemma and the minimal row identity, then the factor-$2$
regrouping estimate (11).  If these three points are accepted, the remaining
argument is a direct compactness-and-amplification contradiction.

\begin{thebibliography}{9}
\bibitem{Braga2021}
B. M. Braga,
\emph{Towards a theory of coarse geometry of operator spaces},
arXiv:2106.15022; Israel J. Math. 259 (2024), 527--558.

\bibitem{BragaOikhberg2022}
B. M. Braga and T. Oikhberg,
\emph{Coarse geometry of operator spaces and complete isomorphic embeddings
into $\ell_1$ and $c_0$-sums of operator spaces},
arXiv:2211.11854; Math. Z. 304 (2023), Article 52.

\bibitem{Pisier}
G. Pisier,
\emph{Introduction to Operator Space Theory},
Cambridge University Press, 2003.

\bibitem{AlbiacKalton}
F. Albiac and N. J. Kalton,
\emph{Topics in Banach Space Theory}, 2nd ed., Springer, 2016.

\bibitem{Ledoux}
M. Ledoux,
\emph{The Concentration of Measure Phenomenon},
American Mathematical Society, 2001.
\end{thebibliography}

\end{document}
